\chapter{The Standard Functions and Their Graphs}

A first course in mathematics for economists assumes that the graphs of the
usual functions are known by sight: which way a cubic runs off at its ends,
where a rational function has a pole, that $\ln x$ never turns around, that
an isoquant is a hyperbola. That knowledge is not decoration either. Every
optimisation problem later on begins with a picture, every limit is checked
against one, and a wrong sketch produces a wrong answer faster than any
arithmetic slip.

This chapter is the catalogue, and the drill that makes it stick. First the
standard shapes; then the features that can be read straight off a formula
without any calculus: intercepts, roots, symmetry, poles and asymptotes,
and the behaviour at the far ends of the axis. The last block leaves the
world of $y=f(x)$ and sketches curves given by an \emph{equation} in $x$
and $y$: circles, ellipses, hyperbolas, budget lines and indifference
curves.

\section{Core Rules}

\begin{rulebox}{The power functions}
Whole-number powers alternate: even powers are symmetric about the
$y$-axis and non-negative, odd powers are symmetric about the origin and
run from $-\infty$ to $+\infty$. Negative powers have a pole at $0$; the
root $\sqrt{x}$ is the upper half of a sideways parabola.
\end{rulebox}

\begin{gallery}
\begin{tikzpicture}
  \begin{axis}[fnaxis, xmin=-2.6, xmax=2.6, ymin=-2.6, ymax=2.6,
      xtick={-2,2}, ytick={-2,2}, title={$y=x$}]
    \addplot[domain=-2.4:2.4]{x};
  \end{axis}
\end{tikzpicture}\hfill
\begin{tikzpicture}
  \begin{axis}[fnaxis, xmin=-2.6, xmax=2.6, ymin=-1.0, ymax=4.4,
      xtick={-2,2}, ytick={2,4}, title={$y=x^{2}$}]
    \addplot[domain=-2.1:2.1]{x^2};
  \end{axis}
\end{tikzpicture}\hfill
\begin{tikzpicture}
  \begin{axis}[fnaxis, xmin=-2.6, xmax=2.6, ymin=-4.4, ymax=4.4,
      xtick={-2,2}, ytick={-4,4}, title={$y=x^{3}$}]
    \addplot[domain=-1.62:1.62]{x^3};
  \end{axis}
\end{tikzpicture}
\end{gallery}

\begin{gallery}
\begin{tikzpicture}
  \begin{axis}[fnaxis, xmin=-0.6, xmax=4.6, ymin=-0.8, ymax=2.6,
      xtick={2,4}, ytick={1,2}, title={$y=\sqrt{x}$}]
    \addplot[domain=0:4.4]{sqrt(x)};
  \end{axis}
\end{tikzpicture}\hfill
\begin{tikzpicture}
  \begin{axis}[fnaxis, xmin=-3.4, xmax=3.4, ymin=-3.4, ymax=3.4,
      xtick={-2,2}, ytick={-2,2}, title={$y=1/x$}]
    \addplot[domain=-3.2:-0.31]{1/x};
    \addplot[domain=0.31:3.2]{1/x};
  \end{axis}
\end{tikzpicture}\hfill
\begin{tikzpicture}
  \begin{axis}[fnaxis, xmin=-3.4, xmax=3.4, ymin=-0.6, ymax=4.4,
      xtick={-2,2}, ytick={2,4}, title={$y=1/x^{2}$}]
    \addplot[domain=-3.2:-0.48]{1/x^2};
    \addplot[domain=0.48:3.2]{1/x^2};
  \end{axis}
\end{tikzpicture}
\end{gallery}

\begin{rulebox}{Exponential, logarithm, absolute value}
$\mathrm{e}^{x}$ is positive everywhere, strictly increasing, and flattens
onto the asymptote $y=0$ as $x\to-\infty$; it passes through $(0,1)$.
$\ln x$ is its mirror image at the line $y=x$: defined only for $x>0$,
strictly increasing, passing through $(1,0)$, with the $y$-axis as a
vertical asymptote. A general exponential $a^{x}$ rises for $a>1$ and falls
for $0<a<1$, always through $(0,1)$; note $\bigl(\tfrac12\bigr)^{x}
=2^{-x}$, so the decreasing case is the increasing one reflected in the
$y$-axis. The absolute value $|x|$ is a V with its vertex at the origin.
\end{rulebox}

\begin{gallery}
\begin{tikzpicture}
  \begin{axis}[fnaxis, xmin=-3.4, xmax=2.4, ymin=-0.8, ymax=4.4,
      xtick={-2,2}, ytick={1,3}, title={$y=\mathrm{e}^{x}$}]
    \addplot[domain=-3.2:1.45]{exp(x)};
  \end{axis}
\end{tikzpicture}\hfill
\begin{tikzpicture}
  \begin{axis}[fnaxis, xmin=-0.8, xmax=5.4, ymin=-3.4, ymax=2.4,
      xtick={1,3,5}, ytick={-2,1}, title={$y=\ln x$}]
    \addplot[domain=0.045:5.2]{ln(x)};
  \end{axis}
\end{tikzpicture}\hfill
\begin{tikzpicture}
  \begin{axis}[fnaxis, xmin=-3.4, xmax=3.4, ymin=-0.8, ymax=3.4,
      xtick={-2,2}, ytick={1,3}, title={$y=|x|$}]
    \addplot[domain=-3:3]{abs(x)};
  \end{axis}
\end{tikzpicture}
\end{gallery}

\begin{rulebox}{What can be read off a formula without calculus}
\begin{itemize}[leftmargin=1.6em,itemsep=0.3ex]
  \item \textbf{$y$-intercept:} evaluate $f(0)$ (if $0\in D$).
  \item \textbf{Roots} \de{(Nullstellen)}: solve $f(x)=0$. For a product,
        set each factor to zero. A factor $(x-a)^{k}$ with $k$ odd makes
        the graph \emph{cross} the axis at $a$; with $k$ even it only
        \emph{touches} it.
  \item \textbf{Sign:} between consecutive roots and poles the sign cannot
        change. Pick one test point per interval; that fixes the whole
        picture.
  \item \textbf{Symmetry:} $f(-x)=f(x)$ gives mirror symmetry in the
        $y$-axis, $f(-x)=-f(x)$ point symmetry in the origin.
  \item \textbf{End behaviour of a polynomial:} only the leading term
        $a_{n}x^{n}$ matters as $x\to\pm\infty$. Even $n$: both ends go the
        same way (up if $a_{n}>0$). Odd $n$: the ends go opposite ways
        (down--left, up--right if $a_{n}>0$).
\end{itemize}
\end{rulebox}

\begin{rulebox}{Rational functions: poles and asymptotes}
Write $f(x)=\dfrac{p(x)}{q(x)}$ with $p,q$ polynomials and
\textbf{cancel common factors first}.
\begin{itemize}[leftmargin=1.6em,itemsep=0.3ex]
  \item A zero of $q$ that does \emph{not} cancel is a \textbf{pole}
        \de{(Polstelle)}: the vertical line $x=a$ is a
        \textbf{vertical asymptote}, and $|f|\to\infty$ there.
  \item A zero of $q$ that \emph{does} cancel is a \textbf{hole}: the graph
        is that of the reduced function, minus the single point $x=a$.
  \item \textbf{Horizontal asymptote}, by comparing degrees:
        $\deg p<\deg q\Rightarrow y=0$; \;$\deg p=\deg q\Rightarrow
        y=\dfrac{\text{leading coeff. of }p}{\text{leading coeff. of }q}$;
        \;$\deg p>\deg q\Rightarrow$ none.
  \item If $\deg p=\deg q+1$, polynomial division gives
        $f(x)=(\text{linear})+\dfrac{\text{remainder}}{q(x)}$, and the
        linear part is an \textbf{oblique asymptote}
        \de{(schiefe Asymptote)}, because the remainder term dies out.
\end{itemize}
\end{rulebox}

\begin{gallery}
\begin{tikzpicture}
  \begin{axis}[fnaxis, xmin=-2.4, xmax=6.4, ymin=-4.4, ymax=6.4,
      xtick={2,4,6}, ytick={-2,2,4}, title={$y=\dfrac{2x+1}{x-3}$}]
    \addplot[domain=-2.2:2.55]{(2*x+1)/(x-3)};
    \addplot[domain=3.42:6.2]{(2*x+1)/(x-3)};
    \addplot[gray, dashed, thin, forget plot, samples=2, domain=-2.2:6.2]{2};
    \addplot[gray, dashed, thin, forget plot, samples=2] coordinates {(3,-4.2)(3,6.2)};
  \end{axis}
\end{tikzpicture}\hfill
\begin{tikzpicture}
  \begin{axis}[fnaxis, xmin=-4.4, xmax=4.4, ymin=-6.4, ymax=6.4,
      xtick={-2,2}, ytick={-4,4}, title={$y=\dfrac{x^{2}+1}{x}$}]
    \addplot[domain=-4.2:-0.17]{(x^2+1)/x};
    \addplot[domain=0.17:4.2]{(x^2+1)/x};
    \addplot[gray, dashed, thin, forget plot, samples=2, domain=-4.2:4.2]{x};
  \end{axis}
\end{tikzpicture}\hfill
\begin{tikzpicture}
  \begin{axis}[fnaxis, xmin=-3, xmax=42, ymin=-3, ymax=32,
      xtick={10,20,30,40}, ytick={5,15,25},
      title={$AC(q)=\dfrac{100+5q}{q}$}]
    \addplot[domain=3.9:40]{(100+5*x)/x};
    \addplot[gray, dashed, thin, forget plot, samples=2, domain=-2.8:40]{5};
  \end{axis}
\end{tikzpicture}
\gcap{Horizontal asymptote $y=2$ and pole $x=3$ \quad\textbullet\quad
oblique asymptote $y=x$ \quad\textbullet\quad average cost falling towards
the asymptote $AC=5$.}
\end{gallery}

\begin{rulebox}{Curves given by an equation in two variables}
Not every curve is a graph $y=f(x)$. The standard ones:
\begin{itemize}[leftmargin=1.6em,itemsep=0.3ex]
  \item \textbf{Line:} $ax+by=c$. Sketch it from the two intercepts
        $x=c/a$ and $y=c/b$; solved for $y$ the slope is $-a/b$. This is
        the budget line $p_{x}x+p_{y}y=M$, with slope $-p_{x}/p_{y}$.
  \item \textbf{Circle:} $(x-x_{0})^{2}+(y-y_{0})^{2}=r^{2}$, centre
        $(x_{0},y_{0})$, radius $r$. A general
        $x^{2}+y^{2}+Dx+Ey+F=0$ is brought to this form by
        \textbf{completing the square} in $x$ and in $y$.
  \item \textbf{Ellipse:} $\dfrac{x^{2}}{a^{2}}+\dfrac{y^{2}}{b^{2}}=1$,
        crossing the axes at $(\pm a,0)$ and $(0,\pm b)$.
  \item \textbf{Hyperbola:} $xy=c$ with $c>0$ lies in the first and third
        quadrant with both axes as asymptotes. Solved for $y$ it is
        $y=c/x$. Every Cobb--Douglas isoquant $K^{\alpha}L^{\beta}=\bar Q$
        and every indifference curve $xy=\bar U$ has this shape.
\end{itemize}
To sketch, either solve for $y$, remembering that a square root splits the
curve into \textbf{two branches} $y=\pm\sqrt{\cdots}$, or read off centre,
radius and intercepts directly.
\end{rulebox}

\section{Practice Problems}

For each polynomial in problems 1--8, give the roots (with multiplicity
where it matters), the $y$-intercept, the symmetry, and the end behaviour
as $x\to\pm\infty$.

\begin{problems}
  \item $f(x)=x^{3}-4x$
  \item $f(x)=x^{4}-5x^{2}+4$
  \item $f(x)=-2x^{3}+6x$
  \item $f(x)=(x-1)^{2}(x+3)$
  \item $f(x)=x^{2}-6x+8$
  \item $f(x)=-x^{2}+4x-4$
  \item $f(x)=x^{5}-x$
  \item A polynomial satisfies $f(x)\to-\infty$ both as $x\to+\infty$ and
        as $x\to-\infty$. What does that say about its degree and its
        leading coefficient?
\end{problems}

In problems 9--20, give all vertical asymptotes, holes, and horizontal or
oblique asymptotes.

\begin{problems}[resume]
  \item $f(x)=\dfrac{1}{x-3}$
  \item $f(x)=\dfrac{2x+1}{x-3}$
  \item $f(x)=\dfrac{x^{2}-1}{x-1}$
  \item $f(x)=\dfrac{3x}{x^{2}-4}$
  \item $f(x)=\dfrac{x^{2}+1}{x}$
  \item $f(x)=\dfrac{x^{2}-4}{x^{2}-9}$
  \item $f(x)=\dfrac{x^{2}+3x+2}{x+1}$
  \item $f(x)=\dfrac{x^{2}-x-6}{x-3}$
  \item $f(x)=\dfrac{2x^{2}+3}{x-1}$
  \item $f(x)=\dfrac{5}{x^{2}+1}$
  \item $AC(q)=\dfrac{100+5q}{q}$ \;(average cost, $q>0$)
  \item $f(x)=\dfrac{x+3}{x^{2}+x-6}$
\end{problems}

\begin{problems}[resume]
  \item For $y=2^{x}$: give the $y$-intercept, the asymptote, the
        monotonicity and the range.
  \item Explain why the graph of $y=\bigl(\tfrac12\bigr)^{x}$ is the graph
        of $y=2^{x}$ reflected in the $y$-axis.
  \item For $y=\ln x$: give the $x$-intercept, the asymptote, the domain
        and the range.
  \item Sketch $y=\mathrm{e}^{-x}$ and state its asymptote and range.
  \item Both $\ln x$ and $\log_{10}x$ pass through $(1,0)$ and rise. What
        exactly distinguishes their graphs?
  \item The adjustment path $y=1-\mathrm{e}^{-x}$, $x\ge0$, is used for
        saturation effects. Give $y(0)$, the asymptote, and say in words
        what the curve does.
\end{problems}

\begin{problems}[resume]
  \item Match each graph (A)--(F) below to one of the formulas
        $\;\ln x$, \;$|x|$, \;$\sqrt{x}$, \;$x^{3}-x$,
        \;$\mathrm{e}^{-x}$, \;$1/x^{2}$.
\end{problems}

\begin{gallery}
\begin{tikzpicture}
  \begin{axis}[fnaxis, xmin=-2.6, xmax=2.6, ymin=-2.2, ymax=2.6,
      xtick={-2,2}, ytick={-2,2}, title={(A)}]
    \addplot[domain=-1.85:1.85]{x^3-x};
  \end{axis}
\end{tikzpicture}\hfill
\begin{tikzpicture}
  \begin{axis}[fnaxis, xmin=-0.7, xmax=4.6, ymin=-0.8, ymax=2.6,
      xtick={2,4}, ytick={1,2}, title={(B)}]
    \addplot[domain=0:4.4]{sqrt(x)};
  \end{axis}
\end{tikzpicture}\hfill
\begin{tikzpicture}
  \begin{axis}[fnaxis, xmin=-3.4, xmax=3.4, ymin=-0.6, ymax=4.4,
      xtick={-2,2}, ytick={2,4}, title={(C)}]
    \addplot[domain=-3.2:-0.48]{1/x^2};
    \addplot[domain=0.48:3.2]{1/x^2};
  \end{axis}
\end{tikzpicture}
\end{gallery}

\begin{gallery}
\begin{tikzpicture}
  \begin{axis}[fnaxis, xmin=-2.4, xmax=3.4, ymin=-0.8, ymax=4.4,
      xtick={-2,2}, ytick={1,3}, title={(D)}]
    \addplot[domain=-1.45:3.2]{exp(-x)};
  \end{axis}
\end{tikzpicture}\hfill
\begin{tikzpicture}
  \begin{axis}[fnaxis, xmin=-3.4, xmax=3.4, ymin=-0.8, ymax=3.4,
      xtick={-2,2}, ytick={1,3}, title={(E)}]
    \addplot[domain=-3:3]{abs(x)};
  \end{axis}
\end{tikzpicture}\hfill
\begin{tikzpicture}
  \begin{axis}[fnaxis, xmin=-0.8, xmax=5.4, ymin=-3.4, ymax=2.4,
      xtick={1,3,5}, ytick={-2,1}, title={(F)}]
    \addplot[domain=0.045:5.2]{ln(x)};
  \end{axis}
\end{tikzpicture}
\end{gallery}

\begin{problems}[resume]
  \item Sketch $f(x)=x^{2}-4x+3$: give the vertex, both intercepts and the
        direction of opening.
  \item Sketch $f(x)=\dfrac{x+1}{x-1}$: give both asymptotes and both
        intercepts.
  \item Sketch $f(x)=\sqrt{x+4}$: give the domain, the range and the
        starting point of the curve.
  \item Sketch $f(x)=|x|-2$: give the vertex and the roots.
  \item Sketch $f(x)=x^{3}-3x^{2}$: give the roots with multiplicity and
        the end behaviour.
  \item Classify as even, odd, or neither: (a) $x^{4}+2$, \;(b) $x^{3}+x$,
        \;(c) $1/x^{2}$, \;(d) $\mathrm{e}^{x}$.
  \item Which of these have a horizontal asymptote, and what is it?
        (a) $\dfrac{3x+1}{x-2}$, \;(b) $\dfrac{x^{2}}{x+1}$,
        \;(c) $\mathrm{e}^{-x}$ as $x\to+\infty$, \;(d) $\ln x$.
  \item On which intervals is $f(x)=\dfrac{(x-1)(x+2)}{x-3}$ positive?
  \item Describe the graph of $f(x)=\sqrt{x^{2}-1}$: domain, symmetry,
        range, and the two points where it meets the $x$-axis.
\end{problems}

In problems 37--45, identify the curve, and sketch it.

\begin{problems}[resume]
  \item $x^{2}+y^{2}=16$. Also solve for $y$ and give the two branches with
        their domains.
  \item $(x-2)^{2}+(y+1)^{2}=9$. Give centre and radius.
  \item $\dfrac{x^{2}}{9}+\dfrac{y^{2}}{4}=1$. Give the four axis
        intercepts.
  \item $xy=6$. Solve for $y$; in which quadrants does the curve lie?
  \item $y^{2}=x$. Is $y$ a function of $x$? Is $x$ a function of $y$?
  \item $2x+3y=12$. Give both intercepts and the slope.
  \item $x^{2}+y^{2}-4x+6y+9=0$. Complete the square; give centre and
        radius.
  \item The isoquant $K^{0.5}L^{0.5}=10$. Solve for $L$ and describe the
        shape for $K>0$.
  \item The budget line $4x+2y=40$ and the indifference curve $xy=50$.
        Give the slope of the budget line and $y$ as a function of $x$ on
        the indifference curve.
\end{problems}

\section{Solutions}

\begin{solutions}
  \item $x^{3}-4x=x(x^{2}-4)=x(x-2)(x+2)$: roots $-2,0,2$, all simple, all
        crossings; $f(0)=0$; odd (only odd powers); degree $3$ with
        positive leading coefficient, so $f\to-\infty$ as $x\to-\infty$ and
        $f\to+\infty$ as $x\to+\infty$.
  \item $x^{4}-5x^{2}+4=(x^{2}-1)(x^{2}-4)$: roots $\pm1,\pm2$;
        $f(0)=4$; even; degree $4$, positive leading coefficient, so both
        ends go to $+\infty$ (a W-shape).
  \item $-2x^{3}+6x=-2x(x^{2}-3)$: roots $0,\pm\sqrt3$; $f(0)=0$; odd;
        degree $3$ with \emph{negative} leading coefficient, so
        $f\to+\infty$ as $x\to-\infty$ and $f\to-\infty$ as $x\to+\infty$.
  \item Roots $x=1$ (double, the graph touches and turns back) and $x=-3$
        (simple, a crossing); $f(0)=(-1)^{2}\cdot3=3$; neither even nor
        odd; degree $3$, leading coefficient $+1$: down on the left, up on
        the right.
  \item $(x-2)(x-4)$: roots $2$ and $4$; $f(0)=8$; vertex at
        $x=3$, $f(3)=-1$; neither even nor odd; parabola opening upwards,
        both ends $+\infty$.
  \item $-(x-2)^{2}$: a double root at $x=2$ (the vertex touches the axis);
        $f(0)=-4$; opens downwards, both ends $-\infty$.
  \item $x^{5}-x=x(x^{4}-1)=x(x^{2}-1)(x^{2}+1)$: roots $-1,0,1$;
        $f(0)=0$; odd; degree $5$, positive leading coefficient: down on
        the left, up on the right.
  \item Both ends behave alike, so the degree is \textbf{even}; both go
        \emph{down}, so the leading coefficient is \textbf{negative}.
  \item Vertical asymptote $x=3$; $\deg p=0<1=\deg q$, so the horizontal
        asymptote is $y=0$. No holes.
  \item Vertical asymptote $x=3$; equal degrees, so
        $y=\tfrac{2}{1}=2$ is the horizontal asymptote.
  \item $\dfrac{(x-1)(x+1)}{x-1}=x+1$ for $x\neq1$: no pole but a
        \textbf{hole} at $(1,2)$. The graph is the line $y=x+1$ with that
        one point punched out; no horizontal asymptote.
  \item $x^{2}-4=(x-2)(x+2)$: poles at $x=\pm2$; $\deg p=1<2=\deg q$, so
        $y=0$.
  \item Pole at $x=0$. Degrees differ by one, so divide:
        $\dfrac{x^{2}+1}{x}=x+\dfrac{1}{x}$, giving the oblique asymptote
        $y=x$.
  \item Poles at $x=\pm3$; equal degrees with leading coefficients $1$ and
        $1$, so the horizontal asymptote is $y=1$.
  \item $\dfrac{(x+1)(x+2)}{x+1}=x+2$ for $x\neq-1$: a hole at $(-1,1)$,
        no pole, no horizontal asymptote.
  \item $\dfrac{(x-3)(x+2)}{x-3}=x+2$ for $x\neq3$: a hole at $(3,5)$.
  \item Pole at $x=1$. Division: $2x^{2}+3=(x-1)(2x+2)+5$, so
        $f(x)=2x+2+\dfrac{5}{x-1}$ and the oblique asymptote is $y=2x+2$.
  \item $x^{2}+1$ never vanishes: no poles at all. $\deg p=0<2$, so
        $y=0$; the graph is a bell-shaped hump with maximum $5$ at $x=0$,
        even.
  \item $AC(q)=\dfrac{100}{q}+5$: pole at $q=0$ and horizontal asymptote
        $AC=5$. Average cost falls towards the marginal cost $5$ as the
        fixed cost $100$ is spread over more units.
  \item $x^{2}+x-6=(x+3)(x-2)$, so
        $f(x)=\dfrac{x+3}{(x+3)(x-2)}=\dfrac{1}{x-2}$ for $x\neq-3$:
        a hole at $\bigl(-3,-\tfrac15\bigr)$, a pole at $x=2$, and the
        horizontal asymptote $y=0$.
  \item $y$-intercept $2^{0}=1$; asymptote $y=0$ (approached as
        $x\to-\infty$); strictly increasing; range $(0,\infty)$.
  \item Because $\bigl(\tfrac12\bigr)^{x}=(2^{-1})^{x}=2^{-x}$: the
        argument $x$ has been replaced by $-x$, and replacing $x$ by $-x$
        mirrors any graph in the $y$-axis.
  \item $x$-intercept $(1,0)$; the $y$-axis $x=0$ is a vertical asymptote
        (as $x\to0^{+}$, $\ln x\to-\infty$); $D=(0,\infty)$, $W=\R$.
  \item The decreasing exponential: it falls from $+\infty$ through
        $(0,1)$ towards the asymptote $y=0$ as $x\to+\infty$;
        $W=(0,\infty)$. It is $\mathrm{e}^{x}$ mirrored in the $y$-axis.
  \item Only the vertical scale: $\log_{10}x=\dfrac{\ln x}{\ln 10}$, so the
        common logarithm is the natural one compressed by the constant
        factor $\ln 10\approx2.303$. Same domain, same asymptote, same
        intercept, same shape.
  \item $y(0)=1-1=0$; as $x\to\infty$, $\mathrm{e}^{-x}\to0$, so the
        asymptote is $y=1$. The curve starts at the origin, rises steeply
        at first and then ever more slowly, approaching the saturation
        level $1$ without ever reaching it.
  \item (A) $x^{3}-x$ \;(odd, three roots); \;(B) $\sqrt{x}$;
        \;(C) $1/x^{2}$ \;(even, positive, pole at $0$);
        \;(D) $\mathrm{e}^{-x}$; \;(E) $|x|$; \;(F) $\ln x$.
  \item $x^{2}-4x+3=(x-1)(x-3)=(x-2)^{2}-1$: opens upwards, vertex
        $(2,-1)$, roots $1$ and $3$, $y$-intercept $3$.
  \item $\dfrac{x+1}{x-1}$: pole $x=1$, horizontal asymptote $y=1$ (equal
        degrees); root $x=-1$; $y$-intercept $f(0)=-1$. Two branches, one
        in each half-plane of the pole.
  \item $D=[-4,\infty)$, $W=[0,\infty)$; the curve starts at $(-4,0)$ and
        rises, ever more slowly: the $\sqrt{\;}$ shape moved four units to
        the left.
  \item A V with vertex $(0,-2)$; roots where $|x|=2$, i.e. $x=\pm2$;
        $y$-intercept $-2$.
  \item $x^{3}-3x^{2}=x^{2}(x-3)$: a double root at $x=0$ (touch) and a
        simple root at $x=3$ (crossing); degree $3$, positive leading
        coefficient: down on the left, up on the right.
  \item (a) even; (b) odd; (c) even, since $1/(-x)^{2}=1/x^{2}$;
        (d) neither.
  \item (a) yes, $y=3$ (equal degrees); (b) no: the numerator degree
        exceeds the denominator degree by one, so there is an oblique
        asymptote $y=x-1$ instead; (c) yes, $y=0$; (d) no: $\ln x$ grows
        without bound, however slowly.
  \item Critical points $-2,1,3$. Test one point per interval:
        $x=-3$ gives $\frac{(-4)(-1)}{-6}<0$; $x=0$ gives
        $\frac{(-1)(2)}{-3}>0$; $x=2$ gives $\frac{(1)(4)}{-1}<0$;
        $x=4$ gives $\frac{(3)(6)}{1}>0$. Hence $f>0$ on
        $(-2,1)\cup(3,\infty)$.
  \item $D=(-\infty,-1]\cup[1,\infty)$ (from $x^{2}-1\ge0$); even, so
        symmetric about the $y$-axis; $W=[0,\infty)$; it meets the axis at
        $(-1,0)$ and $(1,0)$ and rises from there, approaching the lines
        $y=\pm x$ from below for large $|x|$.
  \item Circle, centre $(0,0)$, radius $4$. Solving:
        $y=\pm\sqrt{16-x^{2}}$, two branches, each with domain $[-4,4]$;
        the plus sign is the upper semicircle, the minus sign the lower.
  \item Circle with centre $(2,-1)$ and radius $3$.
  \item Ellipse with $a=3$ and $b=2$: it crosses the axes at
        $(\pm3,0)$ and $(0,\pm2)$.
  \item $y=6/x$, a hyperbola with both axes as asymptotes. Since the
        product is positive, $x$ and $y$ share a sign: the curve lies in
        the first and third quadrants.
  \item A parabola opening to the right. $y$ is \emph{not} a function of
        $x$: the line $x=4$ meets the curve at $y=\pm2$. But $x=y^{2}$
        \emph{is} a function of $y$.
  \item Intercepts $(6,0)$ and $(0,4)$; solved for $y$,
        $y=4-\tfrac23 x$, so the slope is $-\tfrac23$.
  \item $(x^{2}-4x+4)+(y^{2}+6y+9)=-9+4+9$, i.e.
        $(x-2)^{2}+(y+3)^{2}=4$: centre $(2,-3)$, radius $2$.
  \item $K^{0.5}L^{0.5}=10\Rightarrow (KL)^{0.5}=10\Rightarrow KL=100
        \Rightarrow L=\dfrac{100}{K}$. A hyperbola in the first quadrant,
        falling and convex, with both axes as asymptotes: the familiar
        isoquant, along which more capital buys ever less labour saving.
  \item Budget line: $y=20-2x$, slope $-2$ (the price ratio
        $-p_{x}/p_{y}=-4/2$), intercepts $(10,0)$ and $(0,20)$.
        Indifference curve: $y=50/x$, a hyperbola in the first quadrant.
\end{solutions}

\section{Additional Problems}

\begin{problems}
  \item Asymptotes of $f(x)=\dfrac{4x-1}{2x+6}$.
  \item Asymptotes and holes of $f(x)=\dfrac{x^{2}-9}{x+3}$.
  \item Asymptotes and symmetry of $f(x)=\dfrac{x^{3}}{x^{2}-1}$.
  \item Asymptotes and holes of $f(x)=\dfrac{x^{2}-2x}{x^{2}-4}$.
  \item Asymptotes of $f(x)=\dfrac{6}{x^{2}-x-2}$.
  \item Asymptotes of $f(x)=\dfrac{x-4}{x^{2}+4}$.
  \item Roots, symmetry and end behaviour of $f(x)=x^{4}-4x^{2}$.
  \item Roots, sign and end behaviour of $f(x)=(x+2)(x-1)(x-3)$.
  \item Symmetry and asymptote of $f(x)=\dfrac{x}{x^{2}+1}$.
  \item Where do $y=3^{x}$ and $y=3^{-x}$ intersect, and what asymptote do
        they share?
  \item The graphs of $y=\mathrm{e}^{x}$ and $y=\ln x$ are mirror images.
        In which line, and which points correspond to each other?
  \item Domain, range, symmetry and asymptotes of $y=1/x^{2}$.
  \item Domain, symmetry and shape of $y=\sqrt[3]{x}$.
  \item Vertex and roots of $y=|x-2|$.
  \item For the cost function $C(q)=q^{3}-6q^{2}+15q$, tabulate
        $C(0),C(1),\dots,C(5)$ and describe the shape of the curve.
  \item Tabulate $f(x)=x^{2}-2x-3$ for $x=-2,\dots,4$ and sketch it from
        the table; give vertex and roots.
  \item Identify $9x^{2}+4y^{2}=36$ and give its axis intercepts.
  \item Identify $x^{2}-y^{2}=9$ and solve it for $y$.
  \item What curve is $(x+1)^{2}+(y-2)^{2}=0$?
  \item Identify $x^{2}+y^{2}+2x-8=0$: centre and radius.
  \item Compare $y=\sqrt{9-x^{2}}$ with $x^{2}+y^{2}=9$: which is a
        function, and what are its domain and range?
  \item The isocost line $5K+2L=100$: give both intercepts and $L$ as a
        function of $K$.
  \item For the indifference curve $xy=36$, tabulate $y$ at
        $x=1,2,3,4,6,9,12$ and describe the shape.
  \item The demand curve $p=100/q$ has the property that revenue $pq$ is
        constant. Give that constant, the asymptotes, and the economic
        reading of each.
  \item A curve passes through the origin, rises throughout, is concave,
        and flattens onto the horizontal line $y=4$. Which of
        $y=4-\mathrm{e}^{-x}$, \;$y=4(1-\mathrm{e}^{-x})$,
        \;$y=4\mathrm{e}^{-x}$ fits, and why do the other two fail?
\end{problems}

\section{Additional Solutions}

\begin{solutions}
  \item Pole where $2x+6=0$, i.e. $x=-3$; equal degrees, so the horizontal
        asymptote is $y=\tfrac42=2$.
  \item $\dfrac{(x-3)(x+3)}{x+3}=x-3$ for $x\neq-3$: a hole at
        $(-3,-6)$, no pole, no horizontal asymptote (the graph is a line).
  \item Poles at $x=\pm1$. Division: $x^{3}=x(x^{2}-1)+x$, so
        $f(x)=x+\dfrac{x}{x^{2}-1}$ and the oblique asymptote is $y=x$.
        Odd: numerator odd, denominator even.
  \item $\dfrac{x(x-2)}{(x-2)(x+2)}=\dfrac{x}{x+2}$ for $x\neq2$: a hole at
        $\bigl(2,\tfrac12\bigr)$, a pole at $x=-2$, horizontal asymptote
        $y=1$.
  \item $x^{2}-x-2=(x-2)(x+1)$: poles at $x=2$ and $x=-1$; horizontal
        asymptote $y=0$.
  \item $x^{2}+4>0$ always: no poles. Numerator degree smaller, so $y=0$.
  \item $x^{4}-4x^{2}=x^{2}(x^{2}-4)=x^{2}(x-2)(x+2)$: double root at $0$
        (touch), simple roots at $\pm2$; even; degree $4$ with positive
        leading coefficient, so both ends rise to $+\infty$.
  \item Roots $-2,1,3$, all simple. Degree $3$, positive leading
        coefficient: down on the left, up on the right. Signs alternate
        between the roots: negative on $(-\infty,-2)$, positive on
        $(-2,1)$, negative on $(1,3)$, positive on $(3,\infty)$.
  \item Odd (odd numerator over even denominator); no poles since
        $x^{2}+1>0$; horizontal asymptote $y=0$ at both ends. The graph is
        a hump above the axis on the right and its mirror image below on
        the left.
  \item $3^{x}=3^{-x}\iff x=-x\iff x=0$: they meet at $(0,1)$. Both have
        the horizontal asymptote $y=0$, the first as $x\to-\infty$, the
        second as $x\to+\infty$; each is the other reflected in the
        $y$-axis.
  \item In the line $y=x$. The point $(a,\mathrm{e}^{a})$ on the
        exponential corresponds to $(\mathrm{e}^{a},a)$ on the logarithm;
        in particular $(0,1)\leftrightarrow(1,0)$ and
        $(1,\mathrm{e})\leftrightarrow(\mathrm{e},1)$. The two functions
        are inverse to each other.
  \item $D=\R\setminus\{0\}$, $W=(0,\infty)$; even, hence symmetric about
        the $y$-axis; vertical asymptote $x=0$ and horizontal asymptote
        $y=0$. Both branches lie above the axis.
  \item $D=\R$: an odd root accepts negative radicands. It is odd, since
        $\sqrt[3]{-x}=-\sqrt[3]{x}$, strictly increasing, steep near the
        origin and flattening out: the graph of $x^{3}$ mirrored in
        $y=x$.
  \item A V with vertex $(2,0)$; a single (double) root at $x=2$;
        $y$-intercept $|0-2|=2$.
  \item $C(0)=0$, $C(1)=10$, $C(2)=14$, $C(3)=18$, $C(4)=28$, $C(5)=50$.
        Increasing throughout; the increments $10,4,4,10,22$ first fall and
        then rise, so the curve is concave up to about $q=2$ and convex
        afterwards: the S-shaped total cost curve of the textbooks, with
        its inflection at $q=2$.
  \item $f(-2)=5$, $f(-1)=0$, $f(0)=-3$, $f(1)=-4$, $f(2)=-3$, $f(3)=0$,
        $f(4)=5$. The table is symmetric about $x=1$: an upward parabola
        with vertex $(1,-4)$ and roots $-1$ and $3$.
  \item Divide by $36$: $\dfrac{x^{2}}{4}+\dfrac{y^{2}}{9}=1$, an ellipse
        with $a=2$, $b=3$, crossing the axes at $(\pm2,0)$ and
        $(0,\pm3)$.
  \item A hyperbola opening left and right. Solved for $y$:
        $y=\pm\sqrt{x^{2}-9}$, defined for $|x|\ge3$, two branches.
  \item A sum of two squares is zero only if both vanish: the ``curve'' is
        the single point $(-1,2)$.
  \item $(x+1)^{2}-1+y^{2}=8$, so $(x+1)^{2}+y^{2}=9$: a circle with centre
        $(-1,0)$ and radius $3$.
  \item $y=\sqrt{9-x^{2}}$ is a function, the upper semicircle, with
        $D=[-3,3]$ and $W=[0,3]$. The full equation $x^{2}+y^{2}=9$ is not
        a function of $x$: it also contains the lower half.
  \item Intercepts $(20,0)$ on the $K$-axis and $(0,50)$ on the $L$-axis;
        $L=50-2.5K$, a falling line with slope $-2.5$, the ratio of the
        input prices.
  \item $y=36/x$: $36,18,12,9,6,4,3$. A falling, convex hyperbola branch in
        the first quadrant, approaching both axes without touching them:
        halving $x$ doubles $y$.
  \item $pq=100$. The asymptotes are $p=0$ and $q=0$: as the quantity
        grows, the price falls towards zero without reaching it, and no
        finite price clears an arbitrarily small quantity. Revenue stays at
        $100$ everywhere, which is what unit elasticity means.
  \item $y=4(1-\mathrm{e}^{-x})$. It gives $y(0)=0$ and rises to the
        asymptote $4$. \;$y=4-\mathrm{e}^{-x}$ has the right asymptote but
        starts at $y(0)=3$, not at the origin; \;$y=4\mathrm{e}^{-x}$ falls
        from $4$ towards $0$ instead of rising.
\end{solutions}
